\chapter{Dinamometria isoinerziale}

Con la dinamometria isoinerziale si entra nella \textbf{parte specifica della valutazione},
dopo aver trattato i presupposti scientifici che sottendono questo tipo di metodica.

% ---------------------------------------------------------------------------
\section{Le piattaforme di forza}

\textbf{Piattaforma di forza}: strumentazione tra le più importanti e largamente usate, in uso già
dall'inizio degli anni Settanta.

\begin{itemize}
  \item funziona grazie agli \textit{strain gauge} (estensimetri), normalmente \textbf{quattro},
    posizionati agli angoli della piattaforma;
  \item registra le \textbf{variazioni della forza di reazione del terreno in funzione del tempo};
  \item fornisce informazioni \textbf{vettoriali}: le componenti della forza sui tre assi e i
    momenti rispetto al centro di pressione.
\end{itemize}

\rubrica{Campi di applicazione}
\begin{enumerate}
  \item \textbf{locomozione}: variazioni della forza di reazione durante cammino e corsa, in tutte
    le modalità $\rightarrow$ informazioni sui vettori che attraversano tutto il corpo;
  \item \textbf{controllo posturale}: variazione del centro di massa (o centro di gravità) in
    posizione eretta.
\end{enumerate}

\textbf{Sway path}: tracciato delle variazioni nello spazio e nel tempo del centro di massa durante
il mantenimento della stazione eretta; detto anche ``gomitolo'' per la forma del grafico.
\begin{itemize}
  \item è strettamente connesso alla \textbf{capacità di equilibrio} del soggetto;
  \item il grafico fornisce informazioni non solo sul centro del gomitolo --- relativo al centro di
    massa --- ma anche sui tracciati dell'\textbf{arto inferiore destro e sinistro}.
\end{itemize}

% ---------------------------------------------------------------------------
\section{Applicazioni al controllo posturale}

\subsection{Stimolazione plantare ed equilibrio}

Studio condotto con piattaforma di forza per capire come la \textbf{stimolazione dei meccanocettori
plantari} modifichi il controllo posturale (Annino et al., \textit{Somat \& Motor Res}, 2015).

\rubrica{Strumentazione}
Pedana posturale $\rightarrow$ amplificatore $\rightarrow$ personal computer per l'elaborazione dei
tracciati.

\rubrica{Superfici confrontate}
\begin{enumerate}
  \item \textbf{firm}: superficie rigida e liscia;
  \item \textbf{foam}: superficie morbida, come quella delle solette delle scarpe da \textit{running};
  \item \textbf{FTD} (\textit{firm textured disc}): superficie rigida con protuberanze molto piccole,
    in grado di stimolare i meccanocettori plantari; disco di $\phi$ 17\,mm.
\end{enumerate}

\rubrica{Parametri misurati}
\begin{itemize}
  \item \textbf{CoP}: sway path, cioè la distanza coperta dal baricentro durante la durata del test;
  \item \textbf{$V_{A/P}$}: velocità di oscillazione antero-posteriore;
  \item \textbf{$V_{M/L}$}: velocità di oscillazione medio-laterale.
\end{itemize}
Tutte le misure sono state effettuate sia a \textbf{occhi aperti} sia a \textbf{occhi chiusi}.

\rubrica{Risultati}
\begin{itemize}
  \item la superficie stimolante (FTD) ha \textbf{ridotto lo sway path} sia a occhi aperti sia,
    in modo marcato, a occhi chiusi $\rightarrow$ aumento del controllo posturale;
  \item la riduzione ha riguardato anche le \textbf{velocità di spostamento} in entrambe le
    direzioni, antero-posteriore e laterale, in tutte le condizioni;
  \item la \textbf{superficie soffice} è quella che più destabilizza il controllo posturale: non
    stimolando i meccanocettori plantari riduce le afferenze, e il controllo diventa più complicato.
\end{itemize}

\rubrica{Implicazione allenante}
Le esercitazioni di equilibrio avrebbero probabilmente maggiore efficacia se effettuate su
\textbf{superfici rigide e stimolanti} anziché su superfici soffici --- il contrario di come
l'equilibrio viene comunemente allenato, anche nel calcio.

\subsection{Solette testurizzate e fatica muscolare}

Studio sull'effetto delle \textbf{solette testurizzate} sul controllo posturale in posizione eretta
statica dopo affaticamento degli arti inferiori (Palazzo et al., \textit{J Sports Med Phys Fitness},
2017). Le solette contengono piccoli ciottoli non deformabili di varie dimensioni, capaci di
stimolare un numero maggiore di meccanocettori, e vengono inserite nella scarpa.

\rubrica{Metodo}
\begin{itemize}
  \item \textbf{campione}: 10 giovani adulti sani, età media 26,1 $\pm$ 3,07 anni;
  \item \textbf{condizioni sensoriali}: quattro --- con e senza vista, con e senza stimolazione
    plantare naturale;
  \item \textbf{protocollo}: per ogni condizione un singolo trial di 30\,s, in \textit{pre-fatigue}
    e in \textit{post-fatigue};
  \item \textbf{induzione della fatica}: 60\,s di salti continui;
  \item \textbf{misure}: spostamento del centro di pressione (\textit{CoPDISP}), \textit{sway
    velocity}, velocità di oscillazione antero-posteriore e medio-laterale.
\end{itemize}

\rubrica{Risultati e conclusione}
\begin{itemize}
  \item effetto \textbf{stabilizzante} rispetto alle solette di controllo;
  \item riduzione significativa del \textit{CoPDISP} nella condizione a occhi chiusi pre-fatica;
  \item in post-fatica \textbf{tutti} i parametri posturali sono migliorati, sia con sia senza vista;
  \item conclusione: le solette testurizzate a stimolazione rigida migliorano il \textit{CoPDISP}
    \textbf{indipendentemente dalla vista}, fornendo un'informazione sensoriale completa che
    migliora l'equilibrio in condizioni di fatica.
\end{itemize}

\rubrica{Interpretazione}
Buone afferenze propriocettive possono \textbf{ridurre gli effetti negativi della fatica}, tra cui
i livelli di co-contrazione: si rende più efficiente l'azione del sistema neuromuscolare,
soprattutto nel rapporto tra agonisti e antagonisti.

% ---------------------------------------------------------------------------
\section{Analisi del passo su treadmill}

Per l'analisi del passo della corsa si utilizza un \textbf{treadmill multifunzione} dotato di
piattaforma di pressione (\textit{Zebris FDM-TLR}).

\begin{itemize}
  \item \textbf{sensori}: 5376 sensori capacitivi di nuova generazione;
  \item \textbf{frequenza di campionamento}: 100\,Hz;
  \item \textbf{velocità operativa}: 0--14\,km/h.
\end{itemize}

\rubrica{Grandezze analizzate}
\begin{itemize}
  \item \textbf{forza massima} (N) nelle varie fasi del cammino;
  \item \textbf{pressione} (N/cm$^2$): forza esercitata sull'unità di superficie;
  \item \textbf{lunghezza della falcata} (cm);
  \item \textbf{cadenza dei passi} (step/min);
  \item \textbf{doppia fase statica} (\%);
  \item \textbf{swing phase} (\%): percentuale della fase dell'arto oscillante.
\end{itemize}

\rubrica{Ciclo del passo}
$doppio\ appoggio\ iniziale \rightarrow appoggio\ singolo \rightarrow doppio\ appoggio\ terminale
\rightarrow oscillazione$.
\begin{itemize}
  \item doppia fase statica e fase di oscillazione sono caratteristiche del \textbf{cammino};
  \item nella \textbf{corsa} non esiste la doppia fase statica: si hanno solo un \textbf{tempo di
    contatto} e un \textbf{tempo di volo}.
\end{itemize}

% ---------------------------------------------------------------------------
\section{Il principio di misura: la legge di Hooke}

Le \textbf{celle di carico} sono costituite da \textit{strain gauge} e seguono la legge di Hooke.

\textbf{Legge di Hooke}: l'allungamento subìto da un corpo elastico è direttamente proporzionale
alla forza a esso applicata; la costante di proporzionalità è detta \textbf{costante elastica} e
dipende dalla natura del materiale.

\rubrica{Dalla deformazione al segnale}
\begin{enumerate}
  \item il gradiente di forza esercitato sullo \textit{strain gauge} ne provoca la
    \textbf{deformazione};
  \item la deformazione provoca una \textbf{variazione di tensione proporzionale alla forza
    applicata};
  \item calibrando la variazione di tensione con la variazione di forza si ottiene una
    \textbf{relazione lineare} $\rightarrow$ ogni variazione sopra la pedana corrisponde a una
    variazione di forza nell'unità di tempo.
\end{enumerate}

\rubrica{Curva forza-allungamento e limiti di misura}
La curva presenta, nell'ordine: \textbf{regione elastica} $\rightarrow$ \textbf{limite di
proporzionalità} $\rightarrow$ \textbf{limite elastico} $\rightarrow$ \textbf{regione plastica}
$\rightarrow$ \textbf{punto di rottura}.
\begin{itemize}
  \item ogni \textit{strain gauge} ha quindi un proprio \textbf{range di misura};
  \item se la forza applicata supera le capacità elastiche dello strumento, questo \textbf{non
    torna allo stato di quiete} e non fornisce più misure affidabili;
  \item esistono perciò \textit{strain gauge} con scale completamente diverse: quelli delle
    piattaforme, per le variazioni di forza esercitate dall'atleta, e quelli impiegati per misurare
    le tensioni degli smottamenti dei terreni nei terremoti.
\end{itemize}

% ---------------------------------------------------------------------------
\section{L'evoluzione dei test di salto}

Le piattaforme di forza nascono tra la fine degli anni Sessanta e l'inizio degli anni Settanta; i
test iniziali prevedevano invece \textbf{sistemi indiretti}.

\begin{enumerate}
  \item \textbf{Sargent} (1921): misura dell'altezza di salto tramite \textit{reach}, cioè la
    differenza tra il punto toccato dal soggetto su una scala millimetrata a muro e il punto più
    alto raggiunto in volo;
  \item \textbf{Abalakov} (1938): con la fettuccia metrica, differenza tra la posizione eretta e la
    posizione raggiunta nel punto più alto del salto;
  \item \textbf{Davies-Rennie} (1968) e \textbf{Cavagna et al.} (1972): impiego della piattaforma
    di forza;
  \item \textbf{Bosco} (1980): utilizzo del \textbf{tempo di volo} con il tappetino a conduttanza;
  \item \textbf{Vertec} (1990): molto più pratico del test di Sargent, di cui condivide il principio.
\end{enumerate}

\rubrica{Limite dei test indiretti}
Non forniscono \textbf{nessuna informazione sulle capacità neuromuscolari} né sul modo in cui esse
si sviluppano durante il salto: si stabilisce fino a che altezza il soggetto è arrivato, ma non
\textbf{come} l'ha sviluppata.

% ---------------------------------------------------------------------------
\section{Il salto sulla piattaforma di forza}

Con la piattaforma si ottengono informazioni su \textbf{tutti i processi neuromuscolari} che
precedono e caratterizzano la prestazione, non solo sul risultato finale.

\subsection{Lo squat jump}

\textbf{Squat jump}: salto eseguito senza alcun contromovimento $\rightarrow$ modalità
\textbf{concentrica di tipo esplosivo}, in cui il muscolo si accorcia senza essersi allungato
preventivamente.

\rubrica{Lettura del tracciato}
\begin{itemize}
  \item lo sviluppo di forza procede fino al momento dello stacco, seguito dal \textbf{tempo di
    volo} e dalla \textbf{fase di ricaduta};
  \item il \textbf{picco di forza in ricaduta è molto più alto} di quello della fase propulsiva
    $\rightarrow$ spiega perché molti infortuni avvengono in atterraggio più che in spinta;
  \item anche la fase propulsiva resta un momento delicato: il sistema neuromuscolare vi è espresso
    sempre al massimo.
\end{itemize}

\rubrica{Dalla forza alle altre grandezze}
\begin{itemize}
  \item la \textbf{velocità} diventa massima quando la forza diventa minima: forza e velocità stanno
    in rapporto quasi inverso, così come velocità e accelerazione, perché $F = m \cdot a$;
  \item il \textbf{picco di velocità} si raggiunge allo stacco; poi la velocità decade linearmente
    per effetto della gravità, fino all'impatto al suolo --- dove il soggetto sperimenta la forza di
    reazione più alta;
  \item l'\textbf{altezza massima} corrisponde al momento in cui la velocità torna a $0$, cioè il
    punto più alto della parabola, che coincide con la \textbf{metà del tempo di volo}.
\end{itemize}

\subsection{I tre tipi di salto a confronto}

Tracciati forza-tempo registrati su piattaforma: a parità di tempo di volo, lo sviluppo della forza
è completamente differente.

\begin{enumerate}
  \item \textbf{concentrico esplosivo} (squat jump): la forza è sviluppata da un accorciamento senza
    allungamento preventivo;
  \item \textbf{con contromovimento} (CMJ): si osserva una \textbf{fase eccentrica} caratterizzata
    da un peso apparente ridotto rispetto al peso corporeo --- la forza di reazione del terreno
    diminuisce perché il soggetto va verso il basso --- seguita da un'impennata fino al
    \textbf{picco di forza all'inizio della fase concentrica}, che poi diminuisce fino al volo;
  \item \textbf{pliometrico}: l'espressione di forza è caratterizzata da un picco che si esprime
    durante la \textbf{sola fase di contatto}, in cui si distinguono ugualmente una fase eccentrica
    e una concentrica.
\end{enumerate}

\rubrica{Naturalezza del gesto}
Il movimento più naturale è quello determinato dal \textbf{ciclo stiramento-accorciamento}; il
movimento puramente concentrico è meno naturale e difficile da realizzare, se non appreso
precedentemente.

\subsection{Le fasi del salto e le grandezze calcolate}

Sequenza delle fasi rilevate dalla piattaforma: $standing \rightarrow eccentrica \rightarrow
concentrica \rightarrow volo \rightarrow atterraggio \rightarrow standing\ still$; per ciascuna si
seguono \textbf{forza verticale}, \textbf{velocità}, \textbf{altezza di salto} e \textbf{potenza}.

\begin{itemize}
  \item nella \textbf{fase eccentrica} la forza registrata è più bassa del peso corporeo;
  \item il \textbf{picco di forza} si osserva quando la velocità è $0$, nel punto più basso, cioè
    nella posizione di accosciata $\rightarrow$ cade all'\textbf{inizio della fase propulsiva}, non
    alla fine;
  \item di conseguenza, nelle prestazioni con ciclo stiramento-accorciamento la capacità di
    \textbf{reclutare unità motorie ed esprimere forza deve avvenire all'inizio del movimento};
  \item raggiunto il picco la forza diminuisce, perché diminuisce l'inerzia dell'oggetto da muovere
    --- il peso del corpo --- mentre la \textbf{velocità cresce}; dal momento in cui la velocità
    raggiunge il proprio picco interviene la forza di gravità;
  \item la velocità così espressa è \textbf{altamente correlata con l'altezza di salto}, raggiunta a
    metà del tempo di volo.
\end{itemize}

\subsection{Applicazione al sollevamento pesi}

Confronto delle curve forza-tempo di \textbf{strappo} e \textbf{girata} a parità di carico.
\begin{itemize}
  \item la parte iniziale di \textbf{forza esplosiva} è sviluppata nella fase di spinta;
  \item il \textbf{tempo di volo} corrisponde al momento in cui il soggetto ``si infila'' sotto il
    bilanciere, per poi spingerlo;
  \item la \textbf{girata} mostra uno sviluppo di forza maggiore rispetto allo strappo.
\end{itemize}

% ---------------------------------------------------------------------------
\section{Dalla forza alla velocità e allo spostamento}

Il passaggio dalla curva forza-tempo alla curva velocità-tempo, e da questa alla curva
spostamento-tempo, avviene per \textbf{integrazione}:
\[ dF = m\,\frac{dv}{dt} \qquad dF\,dt = m\,dv \qquad \int_{t_1}^{t_2} dF\,dt = m\,\Delta v \]

\rubrica{Procedimento di calcolo numerico}
\begin{enumerate}
  \item si sottrae alla forza la \textbf{costante peso}, che rimane tale $\rightarrow$ si ottiene
    l'\textbf{accelerazione};
  \item si calcola l'area sottesa dalla curva con il \textbf{metodo dei parallelepipedi}: l'area di
    ogni rettangolo vale $a \cdot t$;
  \item la \textbf{sommatoria} delle singole aree dà la velocità $\rightarrow$ $v(t)$;
  \item applicando lo stesso procedimento alla curva $v(t)$ si ottiene lo \textbf{spostamento} in
    funzione del tempo;
  \item da $V(t)$ si ricavano infine $h(t)$ e $P(t)$.
\end{enumerate}

\subsection{Altezza di salto dalla velocità di stacco}

Nota la \textbf{velocità allo stacco} $v_0$ (\textit{take off velocity}), misurata con la pedana di
forza, si applica l'uguaglianza tra energia cinetica ed energia potenziale:
\[ E_c = E_p \qquad \tfrac{1}{2}mv^2 = mgh \qquad \Rightarrow \qquad h_{max} = \frac{v_0^2}{2g} \]
La massa si semplifica tra i due membri.

\subsection{Altezza di salto dal tempo di volo}

Procedimento utilizzato dalle \textbf{pedane conduttive}, che misurano il tempo di volo $t_f$:
\begin{itemize}
  \item \textbf{velocità allo stacco}: $V_{to} = \frac{t_f}{2} \cdot g$ --- accelerazione per tempo,
    quindi già un'integrazione;
  \item \textbf{velocità media}: $V_{media} = \frac{t_f}{4} \cdot g$;
  \item \textbf{altezza}: $h = V_{media} \cdot \frac{t_f}{2} = \frac{t_f}{4} \cdot g \cdot
    \frac{t_f}{2}$.
\end{itemize}
\[ h = \frac{t_f^2 \cdot g}{8} \]
Formula di Asmussen e Bonde Petersen, \textit{Acta Physiol Scand}, 1974.

% ---------------------------------------------------------------------------
\section{Le pedane a contatto}

\subsection{Ergojump}

\textbf{Pedana a conduttanza}: non è altro che un \textbf{interruttore}, costituito da barre di
metallo che vanno in contatto tra loro quando il soggetto è sopra la pedana e si staccano quando è
in aria.
\begin{itemize}
  \item collegata a un \textbf{timer}, permette di ottenere \textbf{tempo di volo} e \textbf{tempo
    di contatto};
  \item esiste anche la versione con \textbf{sistema ottico}, a barre, che sfrutta lo stesso
    principio.
\end{itemize}

\subsection{Wi-Jump}

Sistema che riprende le caratteristiche della pedana a conduttanza rendendola \textbf{wireless}:
$MAT \rightarrow microcontrollore \rightarrow Bluetooth \rightarrow smartphone$ o palmare.
\begin{itemize}
  \item consente di disporre di tutta una serie di parametri \textbf{nell'immediatezza};
  \item è stato \textbf{validato} contro pedana di forza e sistema optoelettronico (Annino et al.,
    2016; 2017);
  \item le relazioni ottenute sono lineari e altamente correlate: $R^2 = 0,975$ con $p < 0,0001$
    ($n = 26$) e $R^2 = 0,96$ con $p < 0,0001$ ($n = 19$) $\rightarrow$ la strumentazione risulta
    \textbf{valida dal punto di vista scientifico}.
\end{itemize}
